\documentclass[11pt]{article}
\usepackage[margin=1in]{geometry}
\usepackage{amsmath,amssymb,amsthm}
\usepackage{graphicx}
\usepackage{hyperref}
\usepackage{xcolor}
\hypersetup{colorlinks=true,linkcolor=blue!55!black,urlcolor=blue!55!black,citecolor=blue!55!black}
\newtheorem{theorem}{Theorem}
\newtheorem{lemma}{Lemma}
\newcommand{\R}{\mathbb{R}}
\newcommand{\perpB}{\mathrel{\perp_B}}
\title{Irregular polygonal Radon planes\\with every admissible number of vertices}
\author{Candidate full solution for arXiv:1809.04760}
\date{11 August 2026}

\begin{document}
\maketitle

\begin{center}
\fbox{\parbox{0.9\linewidth}{\textbf{Status:} full solution likely valid, subject to human review.}}
\end{center}

\section{Source question}

Mandal, Sain, and Paul prove that every Euclidean regular polygon with
$4n+2$ vertices is the unit sphere of a Radon plane and construct irregular
hexagonal examples.  They then ask the following question on page 11 of
arXiv:1809.04760v2~\cite{MSP}.

\begin{quote}
Given any $n\in\mathbb N$, does there exist a Radon plane $X$ such that
$S_X$ is an irregular polygon with $4n+2$ vertices?  If the answer is
affirmative, find a generalized algorithm constructing such planes.
\end{quote}

Here ``regular'' has the source's Euclidean meaning: all side lengths and all
interior angles agree.  We take $\mathbb N=\{1,2,\ldots\}$, as in the source's
polygonal statements.

\begin{figure}[ht]
\centering
\includegraphics[width=\textwidth]{figures/open_problem_crop.png}
\caption{The open question in the original paper, page 11.}
\end{figure}

\section{Proof intuition}

The Radon property belongs to the normed linear structure, so it is invariant
under linear isometry.  Euclidean regularity of a drawn unit polygon is not:
an anisotropic stretch changes its side lengths.  Start with the regular
$(4n+2)$-gon, which the source already proves is Radon, and transport its norm
by the diagonal map $(x,y)\mapsto(\lambda x,y)$.  The transported plane is
isometric to the original Radon plane, but two explicitly chosen edges of its
unit polygon have unequal Euclidean lengths whenever $\lambda\ne1$.

\section{Formal result}

Recall that $x\perpB y$ in a normed space if
$\|x+ty\|\geq\|x\|$ for every $t\in\R$.  A real normed plane is a Radon plane
when this relation is symmetric.

\begin{lemma}[Transport of Birkhoff--James orthogonality]
Let $(V,\|\cdot\|)$ be a real normed space, let $T:V\to V$ be invertible and
define $\|z\|_T=\|T^{-1}z\|$.  Then
\[
 x\perpB y \quad\Longleftrightarrow\quad Tx\perpB Ty
\]
where the relation on the right is computed using $\|\cdot\|_T$.  In
particular, the Radon property is preserved.
\end{lemma}

\begin{proof}
For every real $t$,
\[
 \|Tx+tTy\|_T=\|T^{-1}(Tx+tTy)\|=\|x+ty\|,
 \qquad \|Tx\|_T=\|x\|.
\]
Thus the defining inequalities for the two orthogonality statements are
identical.  Symmetry is therefore transported in both directions.
\end{proof}

\begin{theorem}[Uniform affirmative answer]
For every integer $n\geq1$ and every $\lambda>0$ with $\lambda\ne1$, there is
a Radon norm on $\R^2$ whose unit sphere is an irregular polygon with exactly
$4n+2$ vertices.  More precisely, with $m=4n+2$, let
\[
 v_k=\left(\cos\frac{2\pi k}{m},\sin\frac{2\pi k}{m}\right),
 \quad 0\leq k<m,
 \]
let $P_m=\operatorname{conv}\{v_0,\ldots,v_{m-1}\}$, and set
\[
 T_\lambda(x,y)=(\lambda x,y), \qquad B_{n,\lambda}=T_\lambda(P_m).
\]
Then $B_{n,\lambda}$ is the unit ball of such a norm.
\end{theorem}

\begin{proof}
Since $m$ is even, $P_m$ is centrally symmetric about the origin and has the
origin in its interior, so its Minkowski functional is a norm.  The source
paper proves that this regular $m$-gonal normed plane is Radon~\cite[Theorem
2.7]{MSP}.  The set $B_{n,\lambda}$ is again a centrally symmetric convex body
with the origin in its interior and hence is a unit ball.  Its gauge is
\[
 \|z\|_{n,\lambda}=\|T_\lambda^{-1}z\|_{P_m}.
\]
Consequently $T_\lambda$ is a linear isometry from the regular polygonal norm
onto this new norm.  The lemma shows that the new plane is Radon.  Because
$T_\lambda$ is invertible, it maps the vertices and edges of $P_m$ bijectively
to those of $B_{n,\lambda}$; hence the latter has exactly $m=4n+2$ vertices.

It remains to show that this polygon is not Euclidean regular.  Put
$a=\pi/m$.  The $k$th original edge vector satisfies
\[
 v_{k+1}-v_k
 =2\sin a\left(-\sin((2k+1)a),\ \cos((2k+1)a)\right),
\]
where indices are taken modulo $m$.  Its image therefore has Euclidean length
\[
 L_k=2\sin a\sqrt{\lambda^2\sin^2((2k+1)a)
                         +\cos^2((2k+1)a)}.
\]
For $k=n$, the identity $(2n+1)a=\pi/2$ gives
\[
 L_n=2\lambda\sin a,
\]
whereas
\[
 L_0=2\sin a\sqrt{\lambda^2\sin^2a+\cos^2a}.
\]
If $L_0=L_n$, squaring and cancelling $4\sin^2a$ would yield
\[
 (1-\lambda^2)\cos^2a=0.
\]
But $m\geq6$, so $0<a\leq\pi/6$ and $\cos a\ne0$, while
$\lambda\ne1$.  This is impossible.  Thus two sides have different lengths,
so $B_{n,\lambda}$ is irregular in the source's sense.
\end{proof}

\section{Construction algorithm}

Given $n\geq1$, set $m=4n+2$ and output the vertices
\[
 \boxed{\quad
 w_k=\left(2\cos\frac{2\pi k}{m},\ \sin\frac{2\pi k}{m}\right),
 \quad k=0,\ldots,m-1.\quad}
\]
Their convex hull is the unit ball.  This is the preceding theorem with the
fixed choice $\lambda=2$, so no search or case distinction is needed.

\section{Verification notes}

The argument is symbolic and has no numerical dependency.  The accompanying
regression script is stored at
\begin{center}
\texttt{code/check\_edge\_lengths.py}.
\end{center}
It evaluates all image edge lengths for $n=1,\ldots,100$ at $\lambda=2$ and
checks both the displayed formula and $L_0\ne L_n$.  It is a regression test
only, not part of the proof.

The main human-review points are: (i) the source theorem for the regular
$(4n+2)$-gon; (ii) invariance under the transported norm; and (iii) that the
source uses ``irregular'' in the Euclidean, not affine, sense.  No stronger
claim about classification of all irregular polygonal Radon planes is made.

\section{Bounded novelty check}

The search, performed through 11 August 2026, covered the run's registry,
solution, attempt, and proof-gap indexes; exact-title and exact-question web
searches; arXiv queries for polygonal Radon planes; and combinations of the
phrases \emph{irregular}, \emph{$4n+2$}, \emph{affine image}, and
\emph{affinely regular}.  No later paper explicitly answering the exact
all-orders question was found.  Search results do recognize affine-regular
hexagonal Radon spheres, and the isometric-invariance observation is
elementary.  Accordingly, the result is a complete answer to the stated
question, but its novelty confidence is only moderate-to-low pending a
specialist citation check.

\section{Conclusion and review recommendation}

The source question has an affirmative answer for every $n\geq1$, with the
uniform deterministic construction $T_2(P_{4n+2})$.  We recommend promotion as
\texttt{full\_solution\_likely\_valid}, subject to human verification of the
terminology and the bounded novelty search.

\begin{thebibliography}{9}
\bibitem{MSP}
K. Mandal, D. Sain, and K. Paul,
\emph{A complete characterization of Radon planes whose unit spheres are regular polygons},
arXiv:1809.04760v2, 2018.
\url{https://arxiv.org/abs/1809.04760}.
\end{thebibliography}

\end{document}
